% ─────────────────────────────────────────────────────────────
%  03-examen.tex — un examen avec la classe exam
%  À compiler avec pdfLaTeX (réglage par défaut d'Overleaf).
%  Pour imprimer le corrigé : décommenter \printanswers.
% ─────────────────────────────────────────────────────────────
\documentclass[11pt, addpoints]{exam}   % addpoints : total automatique
% \printanswers                          % affiche les solutions

\usepackage[T1]{fontenc}
\usepackage[french]{babel}
\usepackage{amsmath, amssymb}
\usepackage{siunitx}
\sisetup{output-decimal-marker = {,}}
\usepackage[margin=2.2cm]{geometry}

% ── La classe exam en français ───────────────────────────────
\pointpoints{point}{points}
\hqword{Question}
\hpword{Points}
\hsword{Résultat}
\htword{Total}
\renewcommand{\solutiontitle}{\noindent Solution : }

% ── En-tête et pied de page ──────────────────────────────────
\pagestyle{headandfoot}
\firstpageheader{Mathématique, 3\textsuperscript{e} secondaire}
                {}{Examen --- Relation de Pythagore}
\runningheader{}{}{Relation de Pythagore}
\firstpagefooter{}{Page \thepage\ de \numpages}{}
\runningfooter{}{Page \thepage\ de \numpages}{}

\begin{document}

\noindent Nom : \makebox[8cm]{\hrulefill} \hfill Groupe : \makebox[2cm]{\hrulefill}

\vspace{4mm}
\noindent Cet examen compte \numquestions\ questions, pour un total de
\numpoints\ points. La calculatrice est permise.

\begin{center}
  \gradetable[h][questions]
\end{center}

\begin{questions}

\question[2] Lequel de ces triplets peut représenter les mesures des
côtés d'un triangle rectangle ?
\begin{choices}
  \CorrectChoice $3$, $4$ et $5$
  \choice $4$, $5$ et $6$
  \choice $5$, $11$ et $12$
  \choice $6$, $7$ et $9$
\end{choices}

\question Un triangle rectangle a des cathètes de \qty{6}{\centi\metre}
et de \qty{8}{\centi\metre}.
\begin{parts}
  \part[2] Calcule la longueur de l'hypoténuse.
  \begin{solution}
    $c = \sqrt{6^2 + 8^2} = \sqrt{100} = \qty{10}{\centi\metre}$.
  \end{solution}
  \vspace{\stretch{1}}
  \part[2] Calcule l'aire de ce triangle.
  \begin{solution}
    $A = \dfrac{6 \cdot 8}{2} = \qty{24}{\centi\metre\squared}$.
  \end{solution}
  \vspace{\stretch{1}}
\end{parts}

\question[4] Un téléviseur de format $16:9$ a une diagonale de
\qty{32}{po}. Calcule sa largeur, en pouces, au dixième près.
\begin{solution}[3cm]
  La largeur vaut $16k$ et la hauteur $9k$, donc
  $(16k)^2 + (9k)^2 = 32^2$, soit $337k^2 = 1024$ et
  $k \approx \num{1,743}$. La largeur est d'environ \qty{27,9}{po}.
\end{solution}

\question[3] Le triangle dont les côtés mesurent \qty{7}{\centi\metre},
\qty{24}{\centi\metre} et \qty{25}{\centi\metre} est-il rectangle ?
Justifie ta réponse.
\begin{solution}[3cm]
  Oui : $7^2 + 24^2 = 49 + 576 = 625 = 25^2$. Par la réciproque de la
  relation de Pythagore, le triangle est rectangle.
\end{solution}

\end{questions}

\end{document}
